\chapter{Introdução}

\section{Contextualização}

O Pix consolidou-se como infraestrutura relevante do sistema de pagamentos brasileiro ao permitir transferências e pagamentos instantâneos em disponibilidade contínua. Em 2024, foi o instrumento de pagamento com maior crescimento em quantidade de transações no país, com aumento de 52\%, e respondeu por 47\% das transações de pagamento sem uso de espécie no último trimestre do ano \cite{bcb_pagamentos_2024}. A ampliação desse uso torna necessária a identificação célere de comportamentos de risco, sem que operações legítimas sejam indevidamente interrompidas.

A detecção de fraude financeira constitui um problema de classificação particularmente difícil. Os casos confirmados são raros em comparação às operações legítimas, os padrões podem mudar ao longo do tempo e a confirmação do rótulo pode ocorrer depois da transação. Nessas condições, uma solução pode apresentar alta acurácia global e, ainda assim, falhar na classe de maior interesse \cite{abdallah2016survey,dalpozzolo2018realistic}. Métricas orientadas à classe positiva, como precisão, revocação, F1 e área sob a curva precisão--revocação, são, portanto, mais informativas do que a acurácia isolada \cite{saito2015precision}.

O desempenho preditivo, contudo, não é suficiente em um domínio no qual decisões automatizadas precisam ser examinadas por pessoas. Um escore elevado, desacompanhado dos fatores que o influenciaram, oferece pouco apoio à análise e pode induzir confiança excessiva. O método SHAP permite atribuir a cada variável uma contribuição local para a previsão \cite{lundberg2017unified}. Essa explicação matemática descreve o comportamento do classificador, mas não contém, por si só, contexto normativo ou operacional sobre o Pix.

A arquitetura de \emph{Retrieval-Augmented Generation} (RAG) complementa modelos de linguagem com trechos recuperados de uma coleção documental \cite{lewis2020rag}. Neste trabalho, a recuperação foi projetada como complemento, e não substituta, da explicação SHAP: o modelo estima o risco; SHAP identifica os fatores que influenciaram a saída; e o RAG recupera fontes normativas, operacionais e setoriais relacionadas ao contexto brasileiro. A geração textual, por sua vez, deve permanecer condicionada a essas evidências e declarar insuficiência de suporte quando necessário.

O Banco Central instituiu o Mecanismo Especial de Devolução por meio da Resolução BCB nº 103/2021, posteriormente incorporada ao Regulamento do Pix \cite{bcb_resolucao103_2021,bcb_resolucao1_2020}. Esse histórico evidencia uma exigência para a base documental: preservar a identificação, a versão e a vigência das fontes. A mistura de regras de períodos distintos pode resultar em uma resposta linguisticamente plausível, porém documentalmente incorreta.

\section{Problema de pesquisa}

Modelos de aprendizado de máquina podem reconhecer padrões complexos em dados transacionais, mas sua saída costuma ser técnica e insuficiente para justificar o resultado em linguagem acessível. Em sentido oposto, um modelo de linguagem pode produzir uma narrativa convincente sem fidelidade ao classificador ou à fonte recuperada. O problema consiste em integrar previsão, explicabilidade e recuperação documental sem transformar correlação em causalidade, sem inventar o significado de atributos anonimizados e sem ocultar as limitações do conjunto de dados.

O estudo utiliza o conjunto IEEE-CIS Fraud Detection, composto por transações de comércio eletrônico fornecidas pela Vesta \cite{ieee_cis_2019}. O conjunto não contém transações Pix e parte de seus atributos é anonimizada. Consequentemente, ele serve como base para uma prova de conceito de classificação de risco e integração técnica, mas não permite afirmar que o desempenho observado represente uma operação Pix real.

\subsection{Questão de pesquisa}

\textbf{Como integrar um modelo de detecção de fraude, explicações locais por SHAP e recuperação de documentos por RAG para apoiar explicações rastreáveis em português, preservando a fidelidade ao modelo e explicitando as limitações de uma prova de conceito baseada no IEEE-CIS?}

\section{Objetivos}

\subsection{Objetivo geral}

Desenvolver e analisar um protótipo híbrido que combine aprendizado de máquina, SHAP e RAG para classificar risco de fraude em dados transacionais e apoiar a produção de explicações em português fundamentadas em evidências do modelo e em documentos verificáveis relacionados ao contexto Pix brasileiro.

\subsection{Objetivos específicos}

\begin{enumerate}
  \item preparar e caracterizar o conjunto IEEE-CIS, documentando desbalanceamento, valores ausentes, anonimização e diferenças em relação ao domínio Pix;
  \item treinar e comparar XGBoost, Random Forest e uma linha de base por meio de métricas adequadas à classe rara;
  \item aplicar SHAP para identificar, global e localmente, as variáveis que influenciam as previsões;
  \item manter um registro de atributos semânticos que impeça a tradução indevida de colunas anônimas para conceitos Pix;
  \item organizar uma base documental versionada com normas, materiais oficiais e fontes setoriais selecionadas;
  \item implementar a recuperação semântica e preparar a geração de explicações com indicação das fontes utilizadas;
  \item definir um protocolo para avaliar a qualidade da recuperação, a coerência da resposta com SHAP, o suporte documental e a clareza para usuários;
  \item disponibilizar código, artefatos permitidos, decisões metodológicas e instruções de reprodução, respeitando limites de licença e privacidade dos dados.
\end{enumerate}

\section{Justificativa}

A Pesquisa Febraban de Tecnologia Bancária 2024 informa que 73\% dos 22 bancos participantes da questão declararam utilizar inteligência artificial na detecção de fraude e lavagem de dinheiro. O relatório também registra estratégias de detecção e resposta entre as prioridades de cibersegurança. Esses resultados não mensuram especificamente fraude no Pix nem representam todo o sistema financeiro, mas evidenciam a pertinência prática de técnicas de inteligência artificial no setor \cite{febraban2024}.

No plano acadêmico, a relevância da proposta está na integração controlada de três tipos de evidência. O escore preditivo indica quanto a observação se aproxima da classe aprendida pelo modelo; SHAP descreve quais entradas influenciaram essa saída; e o RAG acrescenta contexto documental rastreável. A separação explícita dessas evidências reduz o risco de apresentar uma narrativa do modelo de linguagem como se fosse previsão, causalidade ou conteúdo normativo.

O trabalho também discute limitações frequentemente omitidas em protótipos, entre elas a mudança de distribuição, o vazamento de dados durante a preparação, a seleção inadequada de métricas, a temporalidade das normas e o risco de atribuir significado a variáveis anonimizadas. Esse registro evita caracterizar uma demonstração acadêmica como sistema pronto para decisões financeiras em produção.

\section{Delimitação do trabalho}

O escopo compreende uma prova de conceito executada sobre o IEEE-CIS, com modelos tabulares, explicações SHAP, índice vetorial local e interface demonstrativa. Não fazem parte do escopo:

\begin{itemize}
  \item implantação em uma instituição financeira;
  \item processamento de dados pessoais reais de clientes Pix;
  \item decisão automática de bloqueio, acusação ou devolução;
  \item validação jurídica da elegibilidade ao Mecanismo Especial de Devolução;
  \item comprovação de generalização para transações Pix reais;
  \item comparação exaustiva de todos os modelos de detecção de fraude.
\end{itemize}

Os exemplos destinados à interface devem ser anonimizados ou sintéticos. Os documentos regulatórios precisam ser recuperados com metadados de origem, versão e vigência. A saída é tratada como apoio explicativo sujeito à revisão humana. Enquanto não houver cliente de modelo de linguagem configurado, o protótipo deve declarar a indisponibilidade da geração textual, em vez de simular uma explicação.

\section{Contribuições do trabalho}

As contribuições pretendidas são:

\begin{enumerate}
  \item um pipeline reproduzível de preparação, classificação e explicação local;
  \item um esquema de evidências que conecte SHAP à recuperação documental sem inventar semântica para atributos anônimos;
  \item um corpus documental rastreável, com metadados e controle de integridade;
  \item um protocolo de avaliação separado para desempenho preditivo, recuperação e fidelidade das explicações;
  \item uma interface demonstrativa capaz de expor resultados e indisponibilidades de cada camada;
  \item uma discussão transparente sobre a transferência limitada do IEEE-CIS para o contexto Pix.
\end{enumerate}

\section{Organização do trabalho}

Além desta introdução, o Capítulo 2 apresenta os fundamentos de fraude financeira, Pix, aprendizado de máquina, explicabilidade e RAG. O Capítulo 3 descreve os dados, a preparação, os modelos, o protocolo experimental, o corpus e a integração das três camadas. O Capítulo 4 reúne os resultados preditivos e os resultados disponíveis de explicabilidade e recuperação, com suas restrições. Por fim, o Capítulo 5 sintetiza as conclusões, confronta os objetivos com o que foi efetivamente desenvolvido, registra as limitações e indica trabalhos futuros.
